\documentclass[11pt]{article}

\usepackage[a4paper,margin=2.4cm]{geometry}
\usepackage{amsmath}
\usepackage{booktabs}
\usepackage{enumitem}
\usepackage[hidelinks]{hyperref}
\usepackage{siunitx}

\newcommand{\file}[1]{\texttt{#1}}
\newcommand{\Ox}[1]{{}^{#1}\mathrm{O}}
\newcommand{\Cl}[1]{{}^{#1}\mathrm{C}}
\newcommand{\Ni}[1]{{}^{#1}\mathrm{N}}

\title{Plan for the proton-production\\cross-section comparison}
\author{J.J. G\'omez-Cadenas}
\date{\today}

\begin{document}
\maketitle

\begin{abstract}
This note defines the implementation of a reproducible comparison of
experimental, evaluated and Geant4 positron-emitter production cross sections.
The result will be reported in the independent document
\file{docs/xsections\_comparison.tex}.  The study covers five proton-induced
production channels from \SIrange{5}{120}{MeV}, including a dedicated Geant4
thin-target calculation with the nominal \file{QGSP\_BIC\_HP} physics list.
Every data transformation, calculation, table and figure will be generated by
version-controlled code in the repository.
\end{abstract}

\section{Objective and scientific scope}

The comparison will answer three questions:
\begin{enumerate}[nosep]
  \item How consistent are the available measurements for each production
        channel?
  \item How well do the JENDL and LANL/ENDF evaluations reproduce those
        measurements?
  \item How well does the nominal Geant4 prediction reproduce their energy
        dependence and normalization?
\end{enumerate}

The frozen channel set is
\begin{align}
  \Ox{16}(p,x)\Ox{15},\qquad
  \Cl{12}(p,x)\Cl{11},\qquad
  \Ox{16}(p,x)\Cl{11},\qquad
  \Ni{14}(p,x)\Ni{13},\qquad
  \Ox{16}(p,x)\Ni{13}.
\end{align}
The primary comparison range is \SIrange{5}{120}{MeV}.  Experimental campaigns
remain separate throughout this study.  The comparison reports the data and
model differences directly; construction of recommended curves and uncertainty
replicas follows in the cross-section-systematics implementation.
At the distal end, $\Ox{16}(p,x)\Ni{13}$ is dominated by
$\Ox{16}(p,\alpha)\Ni{13}$.  Its approximately \SI{5.5}{MeV} threshold gives
this channel direct leverage on the end of the activity profile.

\section{Input data}

The current external inputs stored under \file{data/xsections/raw/} are
\begin{itemize}[nosep]
  \item the pinned EXFOR snapshot containing 67 selected experimental series;
  \item the JENDL-4.0/HE residual-production tables;
  \item the LANL/ENDF proton evaluations distributed through TENDL.
\end{itemize}
Two additions complete the five-channel input set: the available
$\Ox{16}(p,x)\Ni{13}$ EXFOR series and the IAEA recommended
$\Ox{16}(p,\alpha)\Ni{13}$ excitation function.  The expanded EXFOR selection
will be pinned to its source commit, and its new series count will be written
into the data provenance and the comparison report.  The IAEA recommendation
enters as a separate evaluated dataset with its tabulated uncertainty.
The Masuda oxygen measurements will be identified as shape information because
their reported curves share an external normalization.  Each experimental
series will preserve its source identifier, energy range, reported energy
uncertainty, reported cross-section uncertainty and normalization metadata.

The simulation input is generated locally: effective residual-production
curves from Geant4 11.4.x using \file{QGSP\_BIC\_HP}.  These curves are simulation
products with complete run metadata rather than imported nuclear-data files.

\section{Repository structure}

The implementation will use the following layout:
\begin{verbatim}
xsections_g4/
|-- CMakeLists.txt
|-- thin_target.cc
|-- include/
|   |-- ThinTargetDetector.hh
|   |-- ThinTargetPrimaryGenerator.hh
|   `-- ResidualRunAction.hh
`-- src/
    |-- ThinTargetDetector.cc
    |-- ThinTargetPrimaryGenerator.cc
    `-- ResidualRunAction.cc

config/
`-- xsections_scan.toml

analysis_transport/xsections/
|-- __init__.py
|-- exfor.py
|-- jendl.py
|-- endf6.py
|-- iaea.py
|-- normalize.py
|-- run_g4_scan.py
|-- g4_cross_sections.py
|-- make_comparison.py
`-- validate_comparison.py

tools/
`-- build_xsections_comparison.py

test/xsections/
|-- test_readers.py
|-- test_g4_cross_sections.py
`-- test_comparison.py
\end{verbatim}

The thin-target executable is a sibling of \file{stageA\_transport}.  Its
geometry, run pattern and output are specific to cross-section measurement,
while its Geant4 version, nominal physics list and isotope definitions match
the treatment simulation.

\section{Geant4 thin-target study}

\subsection{Executable}

The executable \file{xsections\_g4/thin\_target.cc} will accept the target
isotope, proton energy, target areal thickness, number of incident protons,
random seed and output location.  It will construct isotopically pure
$\Cl{12}$, $\Ox{16}$ or $\Ni{14}$ targets and count the requested residual
nuclei directly.  The output from each run will contain aggregate counters and
run metadata:
\begin{itemize}[nosep]
  \item incident-proton count;
  \item total proton inelastic-interaction count;
  \item residual count for all five target--residual channels;
  \item incident and interaction-energy summaries;
  \item target density, physical thickness and nuclear areal density;
  \item Geant4 version, physics list, random seed and run identifier.
\end{itemize}

For a channel $c$, the thin-target estimator is
\begin{equation}
  \sigma_c(E) =
  \frac{N_c(E)}{N_p(E)\,\eta_t}\,10^{27}\ \mathrm{mb},
\end{equation}
where $N_c$ is the residual count, $N_p$ is the incident-proton count and
$\eta_t$ is the target nuclear areal density in \si{cm^{-2}}.  The statistical
uncertainty is derived from the residual count.  The target validation records
the inelastic-interaction probability and the proton energy loss.

\subsection{Energy grid and targets}

The tracked configuration \file{config/xsections\_scan.toml} will be the
authoritative definition of the scan.  Its initial 36-point energy grid is
\begin{align*}
  &5, 5.5, 6, 7, 8, 9, 10, 11,\\
  &12, 14, 16, 18, 20, 22, 24, 26, 28, 30,\\
  &35, 40, 45, 50, 55, 60,\\
  &65, 70, 75, 80, 85, 90, 95, 100, 105, 110, 115, 120
  \quad\text{MeV}.
\end{align*}
The scan uses an energy-dependent areal thickness.  At each target and energy,
the pilot selects the largest thickness for which the mean proton energy loss
is below \SI{0.5}{MeV} and the inelastic-interaction probability is below
$10^{-3}$.  A \SI{5}{mg/cm^2} target provides the reference geometry for the
pilot.  The complete thickness table is stored in
\file{config/xsections\_scan.toml} and copied into the run manifest.

Runs proceed in reproducible batches until every active residual channel on a
target accumulates at least 400 nuclei or the run reaches \num{5e8} incident
protons.  Four hundred residuals correspond to a counting uncertainty of
approximately five percent.  Points that reach the proton cap retain their
exact Poisson interval and a precision flag.  Zero-count points receive a
one-sided 95\% Poisson upper limit.  The run manifest records the stopping
condition and achieved precision.

The reference scan uses the energy-dependent physical targets.  Geant4
cross-section enhancement remains an optional runtime optimization
after the unbiased curves pass the closure tests.

\section{Data conversion and analysis code}

\subsection{Source readers}

The format-specific modules have one responsibility each:
\begin{description}[style=nextline,leftmargin=2.5cm]
  \item[\file{exfor.py}]
    Parse CX4 experimental series and preserve their accession, subentry,
    campaign and reported uncertainties.
  \item[\file{jendl.py}]
    Extract the requested residual-production blocks from the JENDL target
    tables.
  \item[\file{endf6.py}]
    Extract the requested residual products from the LANL/ENDF MF=6, MT=5
    records.
  \item[\file{iaea.py}]
    Convert the recommended $\Ox{16}(p,\alpha)\Ni{13}$ excitation function and
    its tabulated uncertainty while preserving the IAEA evaluation metadata.
\end{description}

The command-line entry point \file{normalize.py} will convert all selected
inputs into the common schema defined by \file{data/xsections/SCHEMA.md}.  One
normalized file will correspond to one source series or evaluated curve.

\subsection{Geant4 scan and cross sections}

The script \file{run\_g4\_scan.py} will read
\file{config/xsections\_scan.toml}, invoke the compiled executable for every
target and energy, and update a run manifest.  Completed runs with matching
configuration metadata can be reused when resuming a scan.

The script \file{g4\_cross\_sections.py} will convert the aggregate counters
into the five effective Geant4 curves.  It will calculate the cross section and
counting uncertainty, attach the simulation metadata and write tables using the
same units and channel identifiers as the external sources.

\subsection{Calculations and plots}

The script \file{make\_comparison.py} will perform all numerical comparisons
and generate all document inputs.  Its products are
\begin{itemize}[nosep]
  \item five absolute cross-section figures, one for each channel;
  \item a source-to-Geant4 ratio panel for each channel;
  \item a source-coverage table;
  \item peak cross sections and peak energies;
  \item normalization-comparison metrics over common energy intervals;
  \item each channel's ratio to the Geant4 total inelastic cross section;
  \item a machine-readable table of the reported metrics.
\end{itemize}
Ratios will be evaluated only where the Geant4 curve exceeds a documented
fraction of its channel maximum.  This suppresses unstable divisions in the
threshold region while retaining the region that contributes to production.
The selected fraction will be stored in the tracked analysis configuration and
printed in the report metadata.

The script \file{validate\_comparison.py} will check units, channel identities,
source coverage, duplicate point identifiers, finite uncertainties, Geant4 run
statistics and consistency between plotted and tabulated quantities.

\section{Generated products}

The compact numerical and graphical products will use this layout:
\begin{verbatim}
data/xsections/normalized/
    one CSV per experimental or evaluated series

data/xsections/g4/geant4-11.4.x/QGSP_BIC_HP/
    run_manifest.csv
    thin_target_counts.csv
    run_meta.csv

data/xsections/models/
    g4_qgsp_bic_hp_<channel>.csv

docs/figures/xsections/
    five comparison figures

docs/generated/xsections/
    coverage_table.tex
    comparison_metrics.tex

docs/xsections_comparison.tex
    standalone comparison report
\end{verbatim}

The compact Geant4 summaries, model curves, figures and generated table
fragments will be versioned with their generating code.  LaTeX auxiliary files
remain build products.

\section{Comparison document}

The report \file{docs/xsections\_comparison.tex} will be a standalone,
compilable article with the following structure:
\begin{enumerate}[nosep]
  \item motivation and channel scope;
  \item experimental and evaluated sources;
  \item data normalization and uncertainty conventions;
  \item Geant4 thin-target method and validation;
  \item comparison plots and metrics for each channel;
  \item channel-by-channel assessment;
  \item conclusions relevant to the later reweighting implementation;
  \item reproducibility information and exact build commands.
\end{enumerate}

The build entry point \file{tools/build\_xsections\_comparison.py} will
regenerate normalized tables, calculations, figures and LaTeX fragments, run
the validation checks and compile the report.  The Geant4 scan remains an
explicit production command because of its computational cost.  The report
builder will verify that the stored scan metadata match the current tracked
configuration.

\section{Tests and acceptance criteria}

The reader tests will use small pinned fixtures representing each supported
source format.  Calculation tests will verify unit conversion, the
thin-target estimator, counting uncertainties, upper limits and ratio masking.
The comparison test will regenerate a deterministic fixture figure and metric
table.

The study is accepted when
\begin{enumerate}[nosep]
  \item every plotted point is traceable to a normalized table and raw source;
  \item every Geant4 point is traceable to a run-manifest entry;
  \item the Geant4 target-thinness criteria are satisfied;
  \item every active Geant4 point reaches 400 residuals or carries the
        documented proton-cap precision flag or upper limit;
  \item all plotted values and document tables are generated by tracked code;
  \item the validation suite passes;
  \item \file{docs/xsections\_comparison.tex} compiles from the repository
        root using the documented command.
\end{enumerate}

\section{Interface to the systematics implementation}

The comparison fixes the channel validity domains used by the later
reweighting.  Let $\mathcal{C}$ be the five supported proton-induced channels
and $\mathcal{D}_c$ the valid energy interval of channel $c$.  An emitter
produced through channel $c_i$ at energy $E_i$ receives
\begin{equation}
  w_i^{(m)} =
  \begin{cases}
    \displaystyle
    \frac{\sigma_{c_i}^{(m)}(E_i)}
         {\sigma_{c_i}^{(\mathrm{G4})}(E_i)},
      & c_i\in\mathcal{C},\ E_i\in\mathcal{D}_{c_i},\\
    1, & \text{for every other classified production route}.
  \end{cases}
\end{equation}
A secondary proton is eligible when its target, residual and interaction
energy identify a supported channel.  Neutron- and ion-induced reactions,
unmodelled target--residual combinations and listed emitters without an
alternative curve retain unit weight.  Provenance completeness and the
reweightable emitter fraction are separate reported counters.

Cross-section-only reweighting varies residual-production branching while
retaining the nominal Geant4 total inelastic attenuation, proton transport and
dose.  The comparison supplies the channel and total-inelastic curves needed by
a tracked diagnostic that bounds the corresponding first-order fluence change.
The full BIC, INCL++ and Bertini calculations define the separate
hadronic-model envelope.

The propagated results distinguish
\begin{align}
  b_{\mathrm{prod}}^{(m)}
    &=R_{50,\mathrm{prod}}^{(m)}-
      R_{50,\mathrm{prod}}^{(\mathrm{nominal})},\\
  b_{\mathrm{reco}}^{(m)}
    &=R_{50,\mathrm{reco}}^{(m)}-
      R_{50,\mathrm{reco}}^{(\mathrm{nominal})}.
\end{align}
The first quantity is available from the emitter-production profile in this
repository.  The second includes acquisition, detector response and
reconstruction downstream.  For the paired replica distribution, the reported
bias is the median of $b_{\mathrm{reco}}$.  The asymmetric one-standard-
deviation interval is defined by the 16th and 84th percentiles.  When one scalar
is required, the adopted cross-section uncertainty is
\begin{equation}
  u_{\mathrm{xs}}=\frac{q_{84}-q_{16}}{2}.
\end{equation}
Individual evaluated-library shifts remain named diagnostics.  Relative to
nominal BIC, the hadronic envelope is the largest absolute reconstructed-edge
shift obtained with INCL++ or Bertini.

\section{Boundary of the comparison study}

This study concludes with the empirical comparison of the five channels using
the pinned measurements, the IAEA recommendation, two evaluated-library
lineages and the nominal Geant4 Binary Cascade calculation.  It supplies the
evidence needed to design the cross-section models used in the next stage.

The subsequent systematics implementation covers recommended cross-section
fits, replica bands, campaign covariance, event reweighting and propagation to
activity profiles and distal-edge displacement.  The INCL++ and Bertini full
transport calculations form the separate hadronic-model envelope.  Additional
reaction channels and additional literature parametrizations enter later when
the first comparison establishes a material need.

\section{Execution sequence}

\begin{enumerate}
  \item Implement and test the four external-source readers.
  \item Extend the EXFOR selection with $\Ox{16}(p,x)\Ni{13}$ and normalize
        the pinned EXFOR, IAEA, JENDL and LANL/ENDF inputs.
  \item Implement and validate the Geant4 thin-target executable.
  \item Run the configured \file{QGSP\_BIC\_HP} target--energy scan.
  \item Convert the aggregate run counters into channel cross sections.
  \item Generate and validate the five comparisons, ratio panels and tables.
  \item Write \file{docs/xsections\_comparison.tex} around the generated
        results.
  \item Rebuild the full result from tracked inputs and compile the standalone
        document.
\end{enumerate}

\end{document}
